% Nouveau EGA V (cote 25): archive cote 25, site page 51 (PDF page 52).
% typescript. Draft compared: none.
% AI-transcribed by Claude Opus 5.5 (claude-opus-5-5, adaptive thinking, effort high), Message Batches API, batch ega25b (msgbatch_01A7n5YuXubUPNK1zsepi8RG), 2026-11-13; not yet checked by eye.
\scanpage{25/51}
\margin{\hand{\unsure{16|46}}}
une quadrique de l'espace dual, relative à la forme duale ...).   Dans le
cas d'une surface non singulière \xout{de} d'un espace projectif, cette situa-%
tion doit cependant être considérée comme exceptionnelle, cf nº suivant.

\inserted{c) Supposons cependant X lisse de dim 2, géom. intègre propre / k, et f une \uline{immersion}. Alors il résulte de 6.1. \unsure{et} 8.8. \inserted{et 8.14.} que si $\underline{Y}^{\mathrm{sing}} \to \check{P}$ n'est pas gén. fini et inséparable, l'ens. des $\{\xi \in \check{P}$ tel que $Y_{\xi}$ soit \deleted{\unsure{géom.}} séparable sur $k(\xi)\}$ avec \uline{au plus 2} comp. irréd. géom. \unsure{a|--} un compl. de codim $\geq 2$}
% NOTE: The repeated circle-and-horizontal-line call marks connect the boxed top addition to this position, before Nous allons. They are graphic insertion signs, not textual o---.

N\unsure{o|g}us allons examiner de plus près   maintenant le cas des surfaces
(le cas des courbes ne se pose évidemment pas du point de vue irréductibi-%
lité des sections hyperplanes \xout{\unsure{t}} ).

(NB  Je m'aperçois avec effroi que la quadrique n'a pas droit, dans
l'emploi que j'ai fait du mot ``\unsure{r|f}églé'', à être appelée réglée !  Cela doit
être en désaccord avec nos pères, et il faudra donc inventer un mot plus
adéquat pour la notion utilisée ici).

\uline{Proposition} 8.18.  Supposons \xout{X géom intègre} k alg clos, X \xout{intègre et}
\deleted{et} intègre (resp. intègre et normale)\inserted{de dimension $d \geq 2$ et}
propre sur k, soit T une partie finie fermée de X telle que  X-T soit
lisse  et f $|$X-T  soit non ramifié.  P\unsure{o|g}ur que l'ensemble des \hand{$\xi \in$} $\overset{\hand{\vee}}{P}$  tels
que  $Y_{\hand{\xi}}$  soit \xout{\unsure{es}} géom irréd   (resp.  géom intègre)  de dimension  1
\unsure{fait}  un complémentaire de  codimension $\hand{\geq}$ 2, il f et s que les deux
conditions suivantes soient satisfaites:

a)  Pour tout x\hand{$\in$}T, il existe une section hyperplane $Y_{\hand{\xi}}$   (\hand{$\xi \in$}P(k))
passant par x, de dimension \over{1}{$d-1$}, et qui soit irréductible.

b) \inserted[typed]{X'=} X-T est ``réglé'' (sic)  pour f, \uline{ou}  il existe  une section hyper
plane  $Y'_{\hand{\xi}}$  (\hand{$\xi \in$} $\overset{\hand{\vee}}{P}$(k) )  de X',  \xout{singulière} qui soit de dimension \over{1}{d}-1,
singulière et  irréductible.

\deleted{Soit  en effet  Z  la  réunion  de  g($\underline{Y}'^{\mathrm{sing}}$), \xout{et} des hyperplans
de P  définis par les  f(x), x  T.  \xout{et}  Alors pour   P - Z,  Y   est
\uline{lisse} de dimension 1, \xout{\illegible{3} tous ses points sont}
\xout{\unsure{géom non singuliers} \illegible{5} au même Y est lisse au dess}
\xout{des points de}  et comme  il est \inserted{géom} connexe par 6.1. , il est géom irréd.}

\margin{\hand{51}}

% NOTE: The interlinear géom in the canceled proof is above est/connexe and is placed there. Its letterforms do not securely distinguish pen from type; the inherited handwritten layer remains tentative.
